\selectlanguage{english}
Grothendieck devotes these pages to the passage from a small amount of generating data to a factorization theory on a large category. The point is not merely to factor an individual map. The construction must vary with that map, remain compatible with enlargement of the generators, and admit a bound which is independent of the particular object being treated. Only then can it contribute to the construction of categories of models and, through them, to the homotopy theories of diagrams which the manuscript seeks to organize.

The chapter brings together two kinds of control. Lifting properties say what a factorization must accomplish; cardinal filtrations say when the construction has accomplished enough and how large its intermediate objects can become. The former belongs to the language of model categories, while the latter continues the systematic treatment of generation and accessibility in SGA~4, Expos\'e~I. Their interaction is the mathematical subject of this part of Chapter~XV. The later discussion of accessible classes is therefore not a digression. It addresses the problem of replacing a class of prospective generators by a set to which the preceding construction can be applied.

\medskip\noindent\textbf{Attaching solutions and taking a sufficiently long limit.}
For a set $C_0$ of arrows and a map $f:X\to Y$, the first stage attaches, simultaneously, one solution for every square from an element of $C_0$ to $f$. Grothendieck realizes this simultaneous attachment as the filtered colimit of the finite amalgamations. This permits him to work with the relevant pushouts and filtered colimits, without assuming arbitrary coproducts at the outset. The residual map to $Y$ is then subjected to the same operation, and limit stages are formed by filtered colimits. The result is the classical small object construction, expressed as a transfinite sequence of functorial factorizations; compare \citet[Chapter~12]{Riehl2014}.

The decisive step is a factorization of a map from the domain of a generator through an earlier stage. If these domains are $\pi$-accessible in the manuscript's convention, a limit ordinal of cofinality greater than $\pi$ suffices. A square reaching the limit then already occurs at some stage and receives its solution at the next. The word \emph{limit} is necessary: an ordinal with a last predecessor can satisfy the order-theoretic definition of ``grand devant $\pi$'' while providing no further stage in which to attach the required solution. The printed statements omit this qualification. Their usual successor-cardinal choice of stopping ordinal has the required property and repairs them.

The cellular left class is compared with the left lifting class by the familiar retract argument. Factoring a map which already has the left lifting property exhibits it as a retract of its cellular factor. The commentary also supplies the functoriality arguments left to the reader, including compatibility with inclusions of generating sets. One should not confuse completion of the lifting task with stabilization of the sequence: for the sole generator $\varnothing\to1$ in sets, the stages have the form $X\amalg(Y\times\alpha)$. They continue growing after the residual map has become surjective. Garner's algebraic refinement treats the additional structure of chosen fillers and its composition laws \citep{Garner2009}; that structure is stronger than the functoriality established here.

\medskip\noindent\textbf{The arithmetic behind a uniform factorization.}
A cardinal filtration presents a large object by sufficiently directed diagrams of smaller ones. The two base cardinals have distinct roles: $\pi_0$ bounds the complexity of these presentations, while $\pi_1$ provides a uniform accessibility bound for the first bounded level. The estimates then depend on a small category of representatives and on its actual arrows. Counting only isomorphism classes in the arrow category loses the automorphisms and cannot bound individual Hom sets. The one-object category of an arbitrarily large small group gives a particularly short counterexample: its arrows have one isomorphism class, although its unique Hom set has the cardinality of the group.

Choose a small full skeleton $K_1$ of the $\pi_1$-level and put
\[
 \chi=\max\bigl(\aleph_0,\pi_1,|\operatorname{Mor}(K_1)|\bigr),
 \qquad \mu=r^{\pi_0},\qquad r,s\geq\chi.
\]
The corrected counting argument gives the intended forms of the bounds: at most $2^\mu$ isomorphism classes of objects and arrows at level $r$, and at most $s^\mu$ maps from an $r$-bounded object to an $s$-bounded one. The number of choices of a diagram must be counted by exponentiation, not by the printed products. With that change the final estimates still follow. The asserted universal sharpness of the two-coloured-set example does not: under GCH its arrow-orbit count at $r=\aleph_1$ is smaller than $2^r$.

The fixed-point device is robust. Taking a sufficiently large $\rho=2^c$ makes $\rho^{\pi_0}=\rho$ and absorbs the repeated attachment bounds. The manuscript's alternative GCH shortcut needs a cofinality condition: GCH does not imply $\rho^{\pi_0}=\rho$ at singular cardinals of cofinality at most $\pi_0$. The commentary separates these two choices and proves a bound retaining the actual number of generators. If their endpoints lie at level $r$, the input endpoints at level $s$, and the construction is followed through $\alpha$, then
\[
 B=\bigl(\max(2,r,s,|C_0|,|\alpha|)\bigr)^{r^{\pi_0}}
\]
bounds every stage, under the stated closure hypotheses. Thus the principal uniform-factorization conclusion survives the correction of its numerical threshold. Low's later analysis identifies regular-cardinal conditions under which factorization functors preserve presentable objects and a small subcategory detects a combinatorial model structure \citep[Theorem~4.4, Theorem~5.9, Proposition~5.12]{Low2016Heart}.

\medskip\noindent\textbf{Recognizing a category of models.}
Section~4 separates the elementary logic of lifting from the smallness needed to produce the factorizations. A precouple supplies lifting and factorization; closing its two classes under retracts gives the corresponding weak factorization system. Two compatible systems, together with the cancellation properties of $W$, provide the model-category identities. This distinction explains both the strength of Proposition~4.2 and the defect in the later footnote which treats unclosed precouples as though they were already closed. A set-theoretic counterexample in the commentary satisfies the displayed equalities but fails retract closure.

The class $\Cof_W$ records preservation of weak equivalences by pushout. An arrow in $W\cap\Cof_W$ need not, under the hypotheses printed at that point, remain in $W$ after every cobase change. A four-object diamond lattice makes the difference explicit. The omitted sufficient hypothesis is already present in Chapter~XIII: every map must factor as a $W$-cofibration followed by a weak equivalence. The commentary proves the resulting identification with couniversal weak equivalences. The eventual recognition theorem, Proposition~4.4.6, supplies this factorization through its first weak factorization system, so its principal conclusion is preserved.

For the existence of the second small generating family, modern accessibility hypotheses give a precise answer. Smith's theorem as published by \citet[Theorem~1.7 and Proposition~1.15]{Beke2000} supplies the relevant solution-set argument in a locally presentable category. The formulation in \citet[Proposition~A.2.6.15, updated online version]{Lurie2009HTT} is especially close to the proposed construction from $W$ and $\Cof_W$. These results resolve the problem in their stated setting; they do not justify deleting hypotheses from an arbitrary cardinally filtered category. Once such a combinatorial model category is obtained, its diagram homotopy categories form a strong derivator \citep[Proposition~1.30]{Groth2013}. This identifies the later destination of the machinery without treating this portion of the manuscript as a proof of every derivator axiom.

\medskip\noindent\textbf{Recovering a class from its bounded part.}
Section~5 replaces presentations by maps with presentations by small subobjects. Images, closure of bounded levels under suitable colimits, and exactness of filtered colimits make the poset of small subobjects into a continuous, sufficiently directed presentation of the object. The hypotheses on images are independent in the direction questioned by the source: a cyclic module may contain a submodule needing arbitrarily many generators. The commentary gives the complete polynomial-ring example, and distinguishes a strict monomorphism from an arbitrary proper monomorphism in the preceding image lemmas.

The definition of an accessible class used here tests \emph{every} perfectly adapted presentation. Recognizing membership from cofinally many good small subobjects is therefore initially only a necessary condition. The first version of Proposition~5.6 misses the information carried by nonmonic presentations; Grothendieck himself corrects it in footnote~21. A continuous diagram of two-point and one-point posets makes the retract obstruction explicit. Closure of the bounded part under small filtered colimits repairs the assertion. The remaining comparison is a closed-cofinal-set argument: two cofinal families meet cofinally when both contain the relevant directed suprema. The supplied proofs include the eventual-equality step in Lemma~5.7 and avoid assuming that the image of a continuous cofinal indexing map is itself closed.

At cardinals where the bounded objects are themselves $\pi$-accessible, these arguments give the four descriptions in the corrected Proposition~5.7, the recovery bijection of Corollary~5.8, and filtered-colimit closure. Proposition~5.9 is printed in wider generality. Its proof invokes a slice argument whose supplied justification uses hypotheses no longer retained. We prove the proposition at the cardinals required for the preceding recognition theorem, but do not certify the unrestricted fixed-rank assertion. This is a stated limit of the present verification, not a designation of an open problem. In the verified range, the resulting full subcategory is accessible and accessibly embedded; it need not be cocomplete, as the nonempty sets show \citep[Definition~2.1]{AdamekRosicky1994}.

The commentary keeps these distinctions visible in its apparatus. The edition remains the authority for the historical text; genuine transcription corrections are applied, while slips visible in the edition are retained and explained. The claim register records the source's proof status separately from the repair. Results developed here carry the marker \ours{} in their commentary blocks and no claim of priority.
